\documentclass[10pt,twocolumn]{article}
\usepackage[T1]{fontenc}
\usepackage[utf8]{inputenc}
\usepackage{lmodern}
\usepackage[margin=1.8cm,top=2.4cm,bottom=2cm,columnsep=0.6cm]{geometry}
\usepackage{fancyhdr}
\usepackage{titlesec}
\usepackage{booktabs}
\usepackage{amsmath,amssymb,amsfonts}
\usepackage{microtype}
\usepackage{xcolor}
\usepackage{mdframed}
\usepackage{parskip}
\usepackage{enumitem}
\usepackage{hyperref}

\definecolor{rulered}{RGB}{180,30,30}
\definecolor{boxgray}{RGB}{245,245,245}
\definecolor{keywordgray}{RGB}{100,100,100}

\pagestyle{fancy}
\fancyhf{}
\fancyhead[L]{\textsc{\small Claude Mag}}
\fancyhead[C]{\small\textit{Machine Learning Research Digest}}
\fancyhead[R]{\small 2026-07-13}
\fancyfoot[C]{\small\thepage}
\renewcommand{\headrulewidth}{0.8pt}

\titleformat{\section}{\normalsize\bfseries\color{rulered}}{}{0pt}{}%
  [\vspace{-4pt}\color{rulered}\rule{\columnwidth}{0.4pt}]
\titlespacing{\section}{0pt}{8pt}{4pt}

\hypersetup{colorlinks=true,urlcolor=rulered,linkcolor=black}

\begin{document}
\setlength{\parindent}{0pt}
\setlength{\parskip}{4pt}

{\small\color{keywordgray}\textbf{Keywords:}
  reinforcement learning \textbullet{}
  vision-language models \textbullet{}
  reward design \textbullet{}
  embodied agents}\par\vspace{4pt}

{\LARGE\bfseries\raggedright
  VLM-AR3L}\par\vspace{4pt}

{\large\itshape\raggedright
  Dual VLM Reward Signals---Absolute and Relative---Unlock
  Open-World Reinforcement Learning Without Reward Engineering}\par\vspace{6pt}

{\small\textbf{By} Kuan-Chen Chen, Winston Chen, Wei-Fang Sun \& Min-Chun Hu
  (National Tsing Hua University; NVIDIA NVAITC)}\par\vspace{2pt}

{\small\color{keywordgray}arXiv:2607.00483 \textbullet{}
  IJCAI 2026}
\par\vspace{2pt}\noindent{\color{rulered}\rule{\columnwidth}{0.6pt}}\vspace{6pt}

Designing reward functions for open-ended reinforcement learning
tasks is the field's most persistent bottleneck. Vision-Language
Models offer a promising oracle---they understand goals expressed
in natural language and can evaluate visual observations---but
single-signal VLM reward methods suffer from either sparse
absolute evaluations or noisy pairwise comparisons.
VLM-AR3L combines both: an absolute reward model scoring current
state quality and a relative reward model measuring inter-frame
progress. On four long-horizon Minecraft tasks from MineDojo,
the dual-signal framework achieves 29--41\% relative improvement
over the strongest single-signal VLM baseline.

\begin{mdframed}[backgroundcolor=boxgray,linecolor=rulered,linewidth=1pt,
    innerleftmargin=6pt,innerrightmargin=6pt,
    innertopmargin=6pt,innerbottommargin=6pt]
\textbf{\small AT A GLANCE}\par\vspace{4pt}
\begin{description}[leftmargin=0pt,labelsep=4pt,itemsep=2pt,parsep=0pt,
    font=\small\bfseries]
  \item[Problem:] Absolute VLM rewards are sparse in long-horizon tasks;
    relative VLM rewards lack goal-distance information. Neither alone
    trains effective agents.
  \item[Why it matters:] Open-ended embodied RL requires dense,
    well-shaped reward signals without human-designed reward functions.
  \item[Method:] Separate absolute (scalar state evaluation) and relative
    (pairwise progress) reward MLPs trained on offline VLM preference
    labels; combined as weighted sum for PPO policy training.
  \item[Result:] 29--41\% relative improvement over best single-signal
    VLM baseline on all four MineDojo Minecraft task categories.
  \item[Evidence:] Ablation confirms both reward branches contribute;
    removing either degrades success rate by 20--30\%.
\end{description}
\end{mdframed}

\section{The Big Picture}

Reinforcement learning in open-ended, visually complex environments
requires reward signals that are dense enough to provide learning signal
throughout each episode and semantically grounded enough to evaluate
progress toward linguistically-expressed goals. Human reward engineering
meets both requirements but is expensive and brittle---reward functions
designed for one task rarely transfer.

Vision-Language Models have transformed multimodal understanding: they
can interpret visual scenes in the context of language instructions, making
them natural candidates for reward oracles. However, VLMs are expensive
to query at every RL training step (millions of environment interactions).
VLM-AR3L addresses this by distilling VLM knowledge into lightweight reward
models queried once, offline, before training begins.

\section{Why Existing Approaches Fall Short}

EUREKA and similar methods generate only \emph{absolute} scalar rewards:
the VLM evaluates each frame independently and assigns a score. In
long-horizon tasks (sequences of hundreds to thousands of steps), this
signal is near-zero for most of the trajectory---the agent receives almost
no feedback until it has nearly achieved the goal. Sparse rewards cause
RL algorithms to explore randomly for most of training.

Preference-based methods (RLHF-style relative rewards) compare pairs of
states and infer binary preferences. These avoid sparsity but introduce
a different pathology: the agent can score well on pairwise comparisons
by finding states that look locally better than their neighbors, without
making global progress. This leads to oscillatory behavior around local
optima.

Prior work that combined absolute and relative rewards did so in domains
with dense symbolic state information (chess, video games with score
counters). VLM-AR3L is the first to derive both signals from visual
observations with natural-language goals using a single VLM.

\section{Core Mechanism}

The framework decomposes the reward signal into two complementary streams:

\textbf{Absolute Reward Model (ARM):} A learned scorer
$f_\phi^{abs}: \mathcal{O} \times \mathcal{G} \to [0,1]$ maps a
current visual observation $x_t$ and language goal $g$ to a scalar
task-achievement estimate. The ARM is stable: even a noisy estimate of
absolute task completion prevents catastrophic reward collapse and
anchors the agent's sense of total progress.

\textbf{Relative Reward Model (RRM):} A learned comparator
$f_\psi^{rel}: \mathcal{O} \times \mathcal{O} \times \mathcal{G} \to [-1,1]$
compares consecutive frames $(x_t, x_{t+1})$ in the context of goal $g$
to infer local progress or regression. The RRM is adaptive: it provides
dense, non-zero feedback throughout the trajectory, detecting each small
step toward the goal.

\section{Technical Formulation}

The combined reward at step $t$:
\[
r_t = \alpha \cdot \sigma\!\left(f_\phi^{abs}(\text{enc}(x_t), \text{enc}(g))\right)
+ (1-\alpha)\cdot\tanh\!\left(f_\psi^{rel}(\text{enc}(x_t),\text{enc}(x_{t+1}),\text{enc}(g))\right)
\]

where $\text{enc}(\cdot)$ denotes a frozen CLIP encoder. The ARM is
trained by minimizing squared error against VLM scalar pseudo-labels:
\[
\mathcal{L}_{abs}(\phi) = \mathbb{E}_{(x,g,\hat{r})}\!
\left[\left(f_\phi^{abs}(\text{enc}(x),\text{enc}(g)) - \hat{r}\right)^2\right]
\]

The RRM is trained via binary cross-entropy against VLM pairwise
preference labels $y \in \{0,1\}$:
\[
\mathcal{L}_{rel}(\psi) = -\mathbb{E}_{(x_t,x_{t+1},g,y)}\!
\left[y\log\sigma(f_\psi^{rel}) + (1-y)\log(1-\sigma(f_\psi^{rel}))\right]
\]

Both networks are 2-layer MLPs applied to frozen CLIP features. The
mixing coefficient $\alpha = 0.4$ is selected by validation.

\section{Learning or Inference Procedure}

\begin{enumerate}[itemsep=2pt,parsep=0pt]
  \item Collect a replay buffer of $(x_t, x_{t+1}, g)$ triples from
    random-policy rollouts in the target environment.
  \item Query VLM (LLaMA-Vision-3.1-8B) offline to generate:
    (a)~absolute pseudo-labels $\hat{r} \in [0,1]$ for each frame;
    (b)~pairwise preference labels $y \in \{0,1\}$ for consecutive pairs.
  \item Train $f_\phi^{abs}$ and $f_\psi^{rel}$ on the VLM labels.
  \item Run PPO with combined reward $r_t = \alpha f_\phi^{abs} + (1-\alpha) f_\psi^{rel}$,
    querying only the lightweight reward MLPs (no VLM calls during RL).
  \item Optionally fine-tune reward models every $10^5$ steps using
    additional VLM labels from the current policy's rollouts.
\end{enumerate}

\section{What the Guarantee Says}

VLM-AR3L does not provide formal convergence guarantees, as these would
require assumptions about the VLM's labeling quality. The paper proves
empirically that: (1)~VLM pseudo-labels for both ARM and RRM are
significantly more accurate than random ($>$65\% agreement with human
labels on held-out frames) and (2)~the combined reward is more dense
(non-zero at $>$80\% of trajectory steps) and better correlated with
episode return (Pearson $\rho = 0.71$) than either signal alone
($\rho = 0.52$ for ARM, $\rho = 0.48$ for RRM).

\section{Experimental Findings}

\begin{center}
\small
\begin{tabular}{lcccc}
\toprule
\textbf{Task} & \textbf{VLM-AR3L} & \textbf{EUREKA} & \textbf{VLM-RM} & \textbf{Random} \\
\midrule
Combat Spider  & 51.4\% & 37.2\% & 41.8\% & 8.3\% \\
Milk Cow       & 62.3\% & 48.1\% & 50.6\% & 12.1\% \\
Shear Sheep    & 57.8\% & 41.0\% & 44.3\% & 9.7\% \\
Hunt Cow       & 48.9\% & 36.8\% & 39.2\% & 7.5\% \\
\bottomrule
\end{tabular}
\end{center}

Task success rates averaged over 100 evaluation episodes.
VLM-AR3L achieves 29--41\% relative improvement over the next best
VLM baseline (VLM-RM) across all four tasks.

\section{Ablations and Interpretation}

\textbf{ARM only:} Removing the RRM and using only absolute rewards
reduces success rates by 22--28\% (average 24\%). Training is slower
to converge due to sparse reward signal early in episodes.

\textbf{RRM only:} Removing the ARM and using only relative rewards
reduces success rates by 18--32\% (average 25\%). Agents converge to
oscillatory policies that score well on pairwise comparisons without
completing tasks.

\textbf{Mixing coefficient $\alpha$:} Performance peaks at $\alpha = 0.4$
(40\% absolute, 60\% relative). Lower $\alpha$ increases reward density
but loses task-completion grounding; higher $\alpha$ restores grounding
but reintroduces sparsity.

\textbf{VLM size:} Replacing LLaMA-Vision-3.1-8B with a 3B model
reduces labeling quality (agreement with humans drops from 67\% to 58\%)
and downstream success rates by approximately 8\%, confirming that larger
VLMs produce better reward supervision.

\section*{Reference}

Kuan-Chen Chen, Winston Chen, Wei-Fang Sun, Min-Chun Hu.
\textbf{VLM-AR3L: Vision-Language Models for Absolute and Relative Rewards
in Reinforcement Learning}.
arXiv:2607.00483, July 2026. IJCAI 2026.
\url{https://arxiv.org/abs/2607.00483}

\end{document}
